\documentclass[10pt,a4paper]{amsart}
\usepackage[margin=19mm]{geometry}
\usepackage{amsmath,amssymb,mathtools}
\usepackage[scr=rsfs,cal=euler]{mathalfa}
\usepackage{xfrac}
\usepackage[unicode,hidelinks,bookmarks=false]{hyperref}
\newcommand{\ie}{i.e.\ }
\newcommand{\cf}{cf.\ }
\newcommand{\resp}{respectively\ }
\newcommand{\Subsection}{\subsection}
\providecommand{\nobreakdash}{\nobreak\hbox{-}}
\setlength{\emergencystretch}{2em}
\multlinegap0pt
\newcommand{\bZ}{\mathbb{Z}} 
\usepackage{amssymb}
\usepackage{xcolor}
\usepackage[shortlabels]{enumitem}
\newtheorem{lem}{Lemma}
\newcommand\abs[1]{\left\lvert #1 \right\rvert}
\newcommand\set[1]{\left\{ #1 \right\}}
\title{A linear upper bound on the $\bZ_p$-Ramsey number of graphs with sufficiently large $2$-packing}
\author{Emily Heath, Andrew Simmons}
\date{Source: arXiv:2605.21817v1 (CC BY 4.0)}
\begin{document}
\maketitle

\noindent\emph{Source and licence.} A linear upper bound on the $\bZ_p$-Ramsey number of graphs with sufficiently large $2$-packing. Version arXiv:2605.21817v1 (CC BY 4.0), distributed by arXiv under the Creative Commons Attribution 4.0 International licence, http://creativecommons.org/licenses/by/4.0/, as recorded in the arXiv licence metadata for this exact version and shown on the abstract page https://arxiv.org/abs/2605.21817v1. Full licence notice: \texttt{licenses/arxiv-2605.21817-CC-BY-4.0.txt}.\par\smallskip

\noindent\textit{Editorial scope.} Counted unit: the article's lem lem:p3s (source0.tex lines 289--292). The article's complete original proof of it is delivered with it (source0.tex lines 293--307, one \texttt{proof} environment of 15 lines). No mathematics is written, repaired or re-derived: the statement and the proof are the article's own lines, in the article's own notation, and no retained line is substituted or re-flowed.  No further source-local context is needed: every cross-reference of every retained line resolves inside the two delivered spans.  Nothing was omitted from any retained span, so the package carries no omission record.  No retained line contains a citation, so the wrapper carries no bibliography and no bibliography entry.  No retained line contains a figure environment, an image-inclusion command or drawing code, so no image is delivered and the package declares no asset dependency.  Editorial formatting only, in addition: an AMS \texttt{amsart} preamble that restates the article's own declarations and macros used by the retained lines, namely the environments \texttt{lem} on separate counters, together with the standard packages those lines need.  The retained environments print in this document as Lemma 1 in order of appearance, whereas the article prints the same items with the numbering of its own full document.  The complete native source of the article as distributed in the arXiv e-print for this exact version is bundled unchanged as \texttt{sources/201088/source0.tex}.
\begin{lem}
    Let $k\geq 2$ be an integer. Let $G$ be an $n$-vertex graph with $\delta(G)\geq 1$ which contains a $2$-packing $S=\set{v_1,\ldots,v_{k-1}}$. Then any $k$-coloring of $K_{n+k-1}$ contains a monochromatic copy of $G$ or at least $k-1$ disjoint, non-monochromatic copies of $P_3$.
    \label{lem:p3s}
\end{lem}

\begin{proof}
    The statement is trivial for $k=2$, so assume that $k\geq 3$. Let $H=K_{n+k-1}$ and let $f:E(H)\to \set{1,\ldots,k}$. Let $\set{a_ib_ic_i}_{i=1}^{l}$ be a collection of disjoint, non-monochromatic copies of $P_3$ in $H$ of maximum size. Assume towards a contradiction that $l\leq k-2$. Observe that $K=H\setminus\bigcup_{i=1}^{l}\set{a_i,b_i,c_i}$ is a monochromatic clique on $n+k-1-3l$ vertices; denote its color by $x$. We show that $K$ is contained in a larger monochromatic subgraph of $H$.
    
    Note that we may assume $f(a_ib_i)=x$ for all $i$. If instead both $f(a_ib_i)\neq x$ and $f(b_ic_i)\neq x$ for some $i$, then we may replace $a_ib_ic_i$ with another non-monochromatic path containing an edge of color $x$ as follows. Consider any distinct vertices $u_1,u_2,u_3\in V(K)$, which must exist since $n\geq 2(k-1)$ and therefore $\abs{V(K)}\geq 3$. 
    Since we must have either $f(a_ib_i)\neq f(b_iu_2)$ or $f(b_ic_i)\neq f(b_iu_2)$, we may assume without loss of generality that $f(a_ib_i)\neq f(b_iu_2)$. 
    It must be the case that at least one of $f(a_iu_1)$, $f(b_iu_2)$, or $f(c_iu_3)$ is equal to $x$, as otherwise it would be possible to obtain a larger collection of non-monochromatic paths by replacing $a_ib_ic_i$ with $a_ib_iu_2$ and $u_1u_3c_i$. Thus, we may replace the path $a_ib_ic_i$ with $u_1a_ib_i$, $u_2b_ia_i$, or $u_3c_ib_i$, respectively, increasing the number of paths in the collection with an edge of color $x$ while preserving the property that the clique outside this collection is still monochromatic in color $x$.

    Observe now that $f(a_iu)=x$ for all $i$ and all $u\in K$. Indeed, if $f(a_iu_1)\neq x$ for some $i$ and $u_1\in K$, then we may again find a larger collection of non-monochromatic paths by considering another pair of vertices $u_2,u_3\in V(K)\backslash\{u_1\}$. In particular, if $f(c_iu_3)=x$, replacing $a_ib_ic_i$ with $u_2u_1a_i$ and $u_3c_ib_i$ yields a larger collection of paths. If $f(c_iu_3)\neq x$, replacing $a_ib_ic_i$ with $b_ia_iu_1$ and $u_2u_3c_i$ yields the same contradiction. 

    Now, we may assume for all $i$ that either (i) $f(a_ic_i)=x$ and $f(c_iu)=x$ for all $u\in K$, or (ii) $f(b_iu)=x$ for all but possibly one vertex in $V(K)$. If not, then there is some $i$ such that both (i) and (ii) fail. If $f(c_iu)=x$ for all $u\in K$ but $f(a_ic_i)\neq x$, then we can replace $a_ib_ic_i$ in our collection of non-monochromatic paths with $b_ia_ic_i$, resulting in a path which satisfies (ii).  Otherwise, there exist distinct $u_1,u_2\in V(K)$ such that $f(c_iu_1)\neq x$ and $f(b_iu_2)\neq x$. A contradiction follows, since choosing $u_3\in V(K)$ distinct from $u_1,u_2$ and replacing $a_ib_ic_i$ with $a_ib_iu_2$ and $u_3u_1c_i$ yields a larger collection of paths. 

    Next, observe that $f(a_ia_j)=x$ for all $i\neq j$. If instead $f(a_ia_j)\neq x$ for some $i\neq j$, then we can find a larger collection of non-monochromatic paths chosen based on whether $a_ib_ic_i$ and $a_jb_jc_j$ satisfy (i) or (ii) above. If $a_ib_ic_i$ and $a_jb_jc_j$ both satisfy (i) above, then for any distinct $u_1,u_2,u_3\in V(K)$, we can use $u_1c_ib_i$, $u_2c_jb_j$, and  $u_3a_ia_j$. If $a_ib_ic_i$ satisfies (i) and $a_jb_jc_j$ satisfies (ii), then we can pick distinct $u_1,u_2,u_3\in V(K)$ so that $f(u_2b_j)=x$ and use $u_1c_ib_i$, $u_2b_jc_j$, and $u_3a_ia_j$. And if both satisfy (ii), then we can pick distinct $u_1,u_2,u_3\in V(K)$ so that $f(u_1b_i)=f(u_2b_j)=x$ and use $u_1b_ic_i$, $u_2b_jc_j$, and $u_3a_ia_j$.
    
    We now describe how to embed a monochromatic copy of $G$ using all of the $a_i$, one of $b_i$ or $c_i$ for each $i$, and all but one vertex in $V(K)$. Indeed, using at least two vertices from each path will be enough since $\abs{V(H)}-l\geq n+1$. Let $b_i'=c_i$ if $a_ib_ic_i$ satisfies (i) above. Otherwise let $b_i'=b_i$, and if it exists, let $u_i$ denote the vertex in $K$ such that $f(b_i'u_i)\neq x$. Note that the $u_i$ are not necessarily distinct. Let $v_1,\ldots,v_{k-1}$ denote the vertices in $S$. For $1\leq i\leq l$, map $v_i$ and one of its neighbors to $b_i'$ and $a_i$, respectively. While there exist $i\neq j$ such that $\deg(v_i),\deg(v_j)\geq 2$ and $u_i$ and $u_j$ exist and are distinct, map an unmapped neighbor of $v_i$ to $u_j$ and an unmapped neighbor of $v_j$ to $u_i$. Afterward, if there exists $i$ such that $\deg(v_i)\geq 2$, $u_i$ exists, and $u_i$ has not yet been assigned a preimage, map $v_{k-1}$ to $u_i$. Embed the rest of $G$ using vertices in $V(K)$, but otherwise arbitrarily. Because the $N[v_i]$ are disjoint, no edge of the form $b_i'u_i$ or $a_ib_j'$ or $b_i'b_j'$ for $i\neq j$ will be used. By our previous observations, the embedding of $G$ is monochromatic in color $x$.
\end{proof}

\end{document}
